%%%%%%%%%% TUNI CUSTOMIZATION %%%%%%%%%%
\usepackage{charter}

\usepackage{qrcode}

% http://labmaster.mi.infn.it/wwwasdoc.web.cern.ch/wwwasdoc/TL8/texmf/doc/latex/base/html/fntguide/node18.html
%% commands to declare specific symbols

% https://tex.stackexchange.com/questions/399859/fake-semibold-bold-without-increasing-the-length-of-the-text

\xdef\maintextfont{Kolumbus_variable-Regular}
\xdef\maintextboldfont{Kolumbus_variable-Regular} 
\xdef\maintextboldfont{Kolumbus_bold_variable-Regular} 
\xdef\mathvecboldfont{Kolumbus_bold_vec-Regular} 
\xdef\mathbigsymbols{Kolumbus_variable_4p5-Regular}

\usepackage{unicode-math} % would this work?

\usepackage{amsmath,stackengine,mathtools} % stacked math highlights
\stackMath


% \setmathfont{latinmodern-math}
% \defaultfontfeatures[\maintextfont]
%     { Path = {./fonts/} ,
%       Extension = .otf }

% setmainfont, Bold gives mathbf
\setmainfont[BoldFont=\mathvecboldfont .otf, 
    BoldItalicFont=\maintextboldfont .otf]
    {\maintextfont .otf}

\setmathfont{latinmodern-math.otf}

\setmathfont[BoldFont=\mathvecboldfont  .otf,
    BoldItalicFont=\maintextboldfont .otf]
    {\maintextfont .otf}

\setsansfont[BoldFont=\maintextboldfont  .otf,
    BoldItalicFont=\maintextboldfont  .otf]
    {\maintextfont .otf}

\setromanfont[BoldFont=\maintextboldfont  .otf,
    BoldItalicFont=\maintextboldfont  .otf]
    {\maintextfont .otf}
    
%\setmonofont[BoldFont=\maintextboldfont  .otf,
%    BoldItalicFont=\maintextboldfont  .otf]
%    {\mathbigsymbols .otf}
    
\usepackage[defaultrm]{mathastext} % must be after \set...font cmds

    
    
%%%%%%%%%% TUNI CUSTOMIZATION %%%%%%%%%%
\usefonttheme{professionalfonts}

\newfontface\bigsymbolsfont[BoldFont=\mathbigsymbols .otf, BoldItalicFont=\mathbigsymbols .otf, ItalicFont=\mathbigsymbols .otf]{\mathbigsymbols .otf}           % font
% \bg[2.7]{ \text{\bigsymbolsfont\{}} \! \!    
%%%%%%%%%% Title specs %%%%%%%%%%
\newfontface\titlefont{\maintextfont .otf}           % font
\newcommand{\titlecolor}{\textcolor{mainViolet}}    % color
\setbeamerfont{title}{size=\Large} %Large                 % size
\newcommand{\inserttitlefontDetails}{\titlefont\titlecolor} % for frametitle-template


%%%%%%%%%% Frametitle specs %%%%%%%%%%
% \newfontface\frametitlefont[BoldFont=\maintextboldfont .otf, BoldItalicFont=\maintextboldfont .otf, ItalicFont=\maintextfont .otf]{\maintextfont .otf}  % font
\newfontface\frametitlefont[BoldFont=\maintextboldfont .otf, BoldItalicFont=\maintextboldfont .otf, ItalicFont=\maintextfont .otf]{\maintextboldfont .otf}  % font
\setbeamercolor{frametitle}{fg=mainViolet}
\setbeamerfont{frametitle}{size=\Large}  % \Large

%%%%%%%%%% Subframetitle specs %%%%%%%%%%
\newfontface\subframetitlefont[BoldFont=\maintextboldfont .otf, BoldItalicFont=\maintextboldfont .otf, ItalicFont=\maintextfont .otf]
{\maintextfont .otf}  % font
% Path = fonts/,

%% Figure captions:
% \xdef\figcapfont{OpenSans}
\xdef\figcapfont{\maintextfont}
\newfontface\figurecaptionfont[BoldFont=\maintextboldfont .otf, BoldItalicFont=\maintextboldfont .otf, ItalicFont=\figcapfont .otf]
{\figcapfont .otf}  % font

%% Main text font:
%\setbeamercolor{framesubtitle}{fg=mainViolet}
\setbeamerfont{framesubtitle}{size=\large} % \normalsize \small        %% customize framesubtitle -command:        
\let\oldframesubtitle\framesubtitle

\newfontface\mainfont[BoldFont=\maintextboldfont .otf, BoldItalicFont=\maintextboldfont .otf, ItalicFont=\figcapfont .otf]
{\figcapfont .otf}  % font

\setbeamerfont{normal text}{size=\large} %Large                 % size
% \setbeamerfont{alerted text}{size=\large,family*={Kolumbus_variable_bold-Regular}}
\setbeamerfont{item}{size=\large}
\setbeamerfont{item projected}{size=\large}
\setbeamerfont{itemize item}{size=\large}


%% horizontal distances should be in ex lengths ??
\usepackage{mfirstuc}
\usepackage{titlecaps}
\Addlcwords{my for an the a is}
 
% the below works for regular kolumbus font: 
\addtobeamertemplate{frametitle}{\vspace{6mm}\hspace{-12mm}
\frametitlefont}
% https://tex.stackexchange.com/questions/120718/capitalising-first-letter-of-each-word-in-section-headings-using-book-class


%%%%% PAGE NUMBERING %%%%%%%%%%%%%%%%%%%%
%% Page numbering to top left: %%%%%%%%%%
\newcommand\framenumYpos{2mm}	% 0 mm
\newcommand\framenumXpos{98mm}% default

% Macros for recognizing the aspect ratio:
\makeatletter
\newcommand{\aspectratio}{Aspect ratio is 4:3 (default)}
\ifthenelse{%
  \lengthtest{\beamer@paperwidth=16cm}}% condition (16cm is the width of the frame with an aspect ratio = 16:9)
    {\renewcommand{\aspectratio}{Aspect ratio is 16:9}% 
    % 16:9 aspect ratio
    \renewcommand\framenumXpos{130.5mm}}
    {}%do nothing
\makeatother
% 4:3 aspect ratio
% \renewcommand\framenumXpos{98mm}

%% SHOW ONLY CURRENT FRAME: %%
\addtobeamertemplate{frametitle}{}{%
\textblockorigin{\framenumXpos}{\framenumYpos}
\begin{flushright}
\begin{textblock*}{30mm}[0,0.5](0cm,0cm) \tikz {\node[mainViolet] at (0.8cm,0.5cm)%  
 {\tiny \hyperlink{\TOCcourse}{\insertframenumber}\hspace*{\fill}};%TSOM_TOC
   }%  \hfill or  \hspace*{\fill}
\end{textblock*}
\end{flushright}}

%\usepackage{ulem}
%\useunder{\uuline}{\bfseries}{\textbf}
%\useunder{\uwave}{\bf}{\textbf}
%\let\oldinsertframetitle\insertframetitle
%\renewcommand{\insertframetitle}{\Uline{\oldinsertframetitle}}

%% 10.6.21 edit, to comment this
%\addtobeamertemplate{frametitle}{\vspace{-8.0mm}}{\vspace{0mm}} 

%% 30.12. edit, to comment this
% \addtobeamertemplate{frametitle}{\vskip-1mm}{\vspace{3mm}} 
\addtobeamertemplate{frametitle}{\vspace{-8.0mm}}{\vspace{0mm}} 


%%%%%%%%%% COLORS %%%%%%%%%%
\newcommand{\red}[1]{{\color{red}#1}}  %% You can create your own 
\newcommand{\tuniblue}[1]{{\color{tuniblue}#1}}  
\newcommand{\tunired}[1]{{\color{tunired}#1}} 
\newcommand{\tuniyellow}[1]{{\color{tuniyellow}#1}} 
 \newcommand{\tuniyellowdarker}[1]{{\color{tuniyellowdarker}#1}} 
 
%%%%%%% EMPH %%%%%%%%%
%\DeclareTextFontCommand{\emph}{\bfseries\em}
%\DeclareTextFontCommand{\emphLink}{\bfseries\em}

% https://stackoverflow.com/questions/1812214/latex-optional-arguments
%\usepackage{xparse}
%\usepackage{mfirstuc}

%\makeatletter
%\newcommand*\define[1]{%
%  %\textit{#1}%
%  \index{#1@\protect\capitalisewords{#1}}%
%   }
%\makeatother
% 
%\DeclareDocumentCommand\EMPH{ m g g}{%
%    {\IfNoValueT{#3}%
%    {\define{#1}}%
%    \IfNoValueF {#3} {\define{#3}}%
%    }%
%    {\IfNoValueT{#2}%
%    {\href{https://en.wikipedia.org/wiki/#1}{\emphLink{#1}}}%
%    \IfNoValueF {#2} {\href{#2}{\emphLink{#1}}}%
%    }%
%}
%%% EMPH withouth indexing:                                
%\DeclareDocumentCommand\EMPHn{ m g g}{%
%    {\IfNoValueT{#3}%
%    {}%
%    \IfNoValueF {#3} {\define{#3}}%
%    }%
%    {\IfNoValueT{#2}%
%    {\href{https://en.wikipedia.org/wiki/#1}{\emphLink{#1}}}%
%    \IfNoValueF {#2} {\href{#2}{\emphLink{#1}}}%
%    }%
%}

%% darn, Beamer class has its own commands for overlay cmds:
%\newcommand<>{⟨command name⟩}[⟨argument number⟩][⟨default optional value⟩]{⟨text⟩}
%\newcommand<>{\makered}[1]{{\color#2{red}#1}}

%\newcommand{\Alert}[2][blue]{\textcolor{#1}{\textbf{#2}}}
%%%%%%% ALERT %%%%%%%%%
\DeclareTextFontCommand{\Alert}{\bfseries\em}

% https://stackoverflow.com/questions/1812214/latex-optional-arguments
%\usepackage{xparse}
%\usepackage{mfirstuc}
 
%% EMPH withouth indexing:                                
\DeclareDocumentCommand\Alert{O{red} m g}{%
    {\IfNoValueT{#3}%
    {\textcolor{#1}{\textbf{#2}}}%
    \IfNoValueF {#3} {\href{#3}{\textcolor{#1}{\textbf{#2}}}}%
    }%
}

%% neat underlines
\usepackage{contour}
\usepackage{ulem}
\renewcommand{\ULdepth}{2.3pt}
\renewcommand{\ULthickness}{1.0pt}%
\contourlength{0.7pt}
%\contourlength{1.0pt}



\newcommand{\Uline}[1]{%
	\renewcommand{\ULthickness}{1.0pt}%
	  \uline{\phantom{#1}}%
	  \llap{\contour[50]{white}{#1}}%
}

%% appearing dashed underline in math mode:
%\newcommand\MDashuline[1][+-]{\bgroup\markoverwith{%
%\appear[#1]{\renewcommand{\ULdepth}{5.0pt}%
%\kern.1em%
%\vtop{\kern\ULdepth \hrule height 1.5pt width .5em}% \hrule height h depth d width w \relax
%\kern.15em
%}}\ULon%
%}

\newcommand\Dashuline{\bgroup\markoverwith{%
\renewcommand{\ULdepth}{0.8pt}%
\kern.1em%
\vtop{\kern\ULdepth \hrule height 1.5pt width .5em}% \hrule height h depth d width w \relax
\kern.15em
}\ULon}

\newcommand\MDashuline[1][+(1)-]{\bgroup\markoverwith{%
\appear[#1]{\renewcommand{\ULdepth}{5.0pt}%
\kern.1em%
\vtop{\kern\ULdepth \hrule height 1.5pt width .5em}% \hrule height h depth d width w \relax
\kern.15em
}}\ULon%
}

\newcommand{\cutline}[3][+(1)-]{%
\textcolor{#2}{\MDashuline[#1]{\textcolor{black}{\appear[+(-1)-]{#3}}}}%
}

\newcommand{\Underlinecomment}[4][+(1)-]{
\renewcommand\useanchorwidth{T}%tightens math to normal-looking eq.
\stackunder[-1.5mm]{\cutline[#1]{#2}{ #3 }}{\scriptstyle\textcolor{#2}{\appear[+(-3)-]{\text{#4}}}}\addtocounter{beamerpauses}{-3}
}
    
%https://tex.stackexchange.com/questions/475646/centering-linebreaks-and-raggedright-in-underbrace-text-in-mathmode

%\markoverwith{\kern.13em
%  \vtop{\kern\ULdepth \hrule width .3em}%
%  \kern.13em}\ULon}

\newcommand{\DUTline}[2][red!50]{%
	\renewcommand{\ULthickness}{1.5pt}%
	  \textcolor{#1}{\Dashuline{\phantom{#2}}}%
	  \llap{\contour[50]{white}{#2}}%
}
  
\newcommand\culine[1][red]{%
\bgroup\markoverwith{\textcolor{#1}{\rule[-0.3ex]{2pt}{1.5pt}}}\ULon
}

\newcommand{\CUTline}[2][red!50]{%
	\renewcommand{\ULthickness}{1.5pt}%
	  \textcolor{#1}{\uline{\phantom{#2}}}%
	  \llap{\contour[48]{white}{#2}}%
}

\newcommand{\UTline}[1]{%
	\renewcommand{\ULthickness}{1.5pt}%
	  \uline{\phantom{#1}}%
	  \llap{\contour{white}{#1}}%
}

% Underlined alert
\DeclareDocumentCommand\UAlert{O{red} m g}{%
    {\IfNoValueT{#3}%
    {\textcolor{#1}{\textbf{\UTline{#2}}}}%
    \IfNoValueF {#3} {\href{#3}{\textcolor{#1}{\textbf{\UTline{#2}}}}}%
    }%
}

% Colored (slightly lighter underlined), Underlined alert
\DeclareDocumentCommand\CUAlert{O{red} m g}{%
    {\IfNoValueT{#3}%
    {\textbf{\CUTline[#1!40]{\textcolor{#1}{#2}}}}%
    \IfNoValueF {#3} {\href{#3}{\textcolor{#1}{\textbf{\CUTline[#1]{#2}}}}}%
    }%
}

% Colored (slightly lighter underdash), Dashed alert
\DeclareDocumentCommand\CDAlert{O{red} m g}{%
    {\IfNoValueT{#3}%
    {\textbf{\DUTline[#1!40]{\textcolor{#1}{#2}}}}%
    \IfNoValueF {#3} {\href{#3}{\textcolor{#1}{\textbf{\DUTline[#1]{#2}}}}}%
    }%
}

%%

